\chapter{Bayesian Estimation: Prior, Posterior, MAP \& Bayes Risk}

While Frequentist statistics treats unknown parameters $\theta$ as fixed, non-random physical constants to be uncovered via repeated sampling, \index{Bayesian Statistics}\textbf{Bayesian Statistics} treats $\theta$ as a random variable endowed with a probability distribution that quantifies subjective degrees of belief and uncertainty. In this chapter, we master \index{Prior Distributions}\textbf{Prior Distributions}, \index{Bayes' Rule for Parameters}\textbf{Bayes' Rule for Parameters}, \index{Conjugate Updating}\textbf{Conjugate Updating}, \index{Maximum A Posteriori (MAP)}\textbf{Maximum A Posteriori (MAP)} estimation, and the decision-theoretic minimization of \index{Bayes Risk}\textbf{Bayes Risk} under distinct loss functions.

\vspace{0.5cm}
\begin{tcolorbox}[colback=primarylight, colframe=primarydark, title={\textbf{What You'll Learn in this Chapter}}]
\begin{itemize}
    \item The Bayesian paradigm: combining prior knowledge $\pi(\theta)$ with data likelihood $L(\theta)$ to form the posterior $f(\theta \mid \mathbf{x})$.
    \item Closed-form conjugate updates: Beta-Binomial, Normal-Normal, and Gamma-Poisson.
    \item Bayesian point estimators: Posterior Mean, Posterior Median, and Posterior Mode (MAP).
    \item Decision theory and loss functions: how Squared Error Loss, Absolute Error Loss, and 0-1 Loss derive unique optimal estimators.
    \item Equivalence of MAP estimation to $L_1$ (Lasso) and $L_2$ (Ridge) regularized machine learning.
\end{itemize}
\end{tcolorbox}

\section{The Bayesian Framework and Bayes' Rule for Parameters}

\begin{intuitionbox}{The Bayesian Compromise: Prior vs. Evidence}
A Frequentist (MLE) relies solely on observed data, which can severely overfit on small samples (e.g. concluding a coin is 100\% rigged after 3 Heads in 3 flips). A Bayesian balances their prior belief (e.g. "most coins are roughly 50\% fair") against empirical evidence. As sample size $n \to \infty$, the likelihood $L(\theta)$ grows exponentially narrower and completely overpowers the prior, causing Bayesian estimates to converge asymptotically to the MLE.
\end{intuitionbox}

\begin{definitionbox}{Prior, Likelihood, Marginal Evidence, and Posterior}
\index{Prior, Likelihood, and Posterior}
Let $\mathbf{x} = (x_1, \dots, x_n)$ be observed sample data and $\theta \in \Theta$ be the parameter vector.
\begin{enumerate}
    \item \textbf{Prior Distribution $\pi(\theta)$:} Probability density quantifying beliefs about $\theta$ prior to data collection.
    \item \textbf{Likelihood Function $f(\mathbf{x} \mid \theta) = L(\theta)$:} Probability density of observing data $\mathbf{x}$ given parameter $\theta$.
    \item \textbf{Marginal Evidence $f(\mathbf{x})$:} The normalizing constant integrating over all possible parameter values:
    $$f(\mathbf{x}) = \int_{\Theta} f(\mathbf{x} \mid \theta) \pi(\theta) \, d\theta$$
    \item \textbf{Posterior Distribution $f(\theta \mid \mathbf{x})$:} The updated conditional density of $\theta$ given empirical evidence:
    $$f(\theta \mid \mathbf{x}) = \frac{f(\mathbf{x} \mid \theta) \pi(\theta)}{f(\mathbf{x})} \propto L(\theta) \cdot \pi(\theta)$$
    $$\text{Posterior} \propto \text{Likelihood} \times \text{Prior}$$
\end{enumerate}
\end{definitionbox}

\begin{lstlisting}[language=Python, caption=Data Science: Conjugate Beta-Binomial Bayesian Updating]
import numpy as np
from scipy.stats import beta

# Prior: Beta(alpha=2, beta=2) - weak prior centered at p = 0.5
a_prior, b_prior = 2, 2

# Observed Data: 18 Heads (Successes) out of 20 Coin Flips
k_successes = 18
n_trials = 20
k_failures = n_trials - k_successes

# Conjugate Update: simply add counts to prior hyperparameters!
a_post = a_prior + k_successes  # 2 + 18 = 20
b_post = b_prior + k_failures   # 2 + 2 = 4

# Point estimates from the posterior
map_estimate = (a_post - 1) / (a_post + b_post - 2)
posterior_mean = a_post / (a_post + b_post)
mle_estimate = k_successes / n_trials

print(f"MLE Estimate:            {mle_estimate:.4f} (Overfits to 0.9000)")
print(f"Bayesian MAP Estimate:   {map_estimate:.4f} (Shrunk towards prior)")
print(f"Bayesian Posterior Mean: {posterior_mean:.4f}")
\end{lstlisting}

\begin{videocard}{IITM BS Lecture: Bayesian Estimation}{https://www.youtube.com/watch?v=iMjDOTIb3-Q}
Official lecture introducing Bayesian philosophy, updating priors to posteriors using Bayes' rule.
\end{videocard}

\begin{videocard}{IITM BS Lecture: Choice of Prior and Examples}{https://www.youtube.com/watch?v=SFlzn-0YgV0}
Official lecture discussing uninformative priors, Beta-Binomial conjugate priors, and Normal-Normal conjugate priors.
\end{videocard}

\section{Point Estimators and Decision-Theoretic Bayes Risk}

To extract a single representative point estimate $\hat{\theta}$ from the continuous posterior curve $f(\theta \mid \mathbf{x})$, we specify a \index{Loss Function}\textbf{Loss Function} $L(\theta, \hat{\theta})$ measuring the penalty incurred when predicting $\hat{\theta}$ while the true parameter is $\theta$. The optimal estimator minimizes the \index{Posterior Expected Loss}\textbf{Posterior Expected Loss}:
$$\hat{\theta}_{\text{Bayes}} = \arg\min_{\hat{\theta}} \E_{\theta|\mathbf{x}}\left[ L(\theta, \hat{\theta}) \right] = \arg\min_{\hat{\theta}} \int_{\Theta} L(\theta, \hat{\theta}) f(\theta \mid \mathbf{x}) \, d\theta$$

\begin{formulabox}{Optimal Estimators Under Canonical Loss Functions}
\begin{enumerate}
    \item \textbf{Squared Error Loss $L(\theta, \hat{\theta}) = (\theta - \hat{\theta})^2 \implies$ Posterior Mean:}
    $$\hat{\theta}_{\text{Mean}} = \E[\theta \mid \mathbf{x}] = \int_{\Theta} \theta f(\theta \mid \mathbf{x}) \, d\theta$$
    \item \textbf{Absolute Error Loss $L(\theta, \hat{\theta}) = |\theta - \hat{\theta}| \implies$ Posterior Median:}
    $$\hat{\theta}_{\text{Median}} = \text{Median of } f(\theta \mid \mathbf{x}) \quad \text{such that } \int_{-\infty}^{\hat{\theta}} f(\theta \mid \mathbf{x}) d\theta = 0.5$$
    \item \textbf{0--1 Classification Loss $L(\theta, \hat{\theta}) = \mathbb{I}(|\theta - \hat{\theta}| > \epsilon) \implies$ MAP (Posterior Mode):}
    $$\hat{\theta}_{\text{MAP}} = \arg\max_{\theta} f(\theta \mid \mathbf{x}) = \arg\max_{\theta} \left[ \ln L(\theta) + \ln \pi(\theta) \right]$$
\end{enumerate}
\end{formulabox}

\begin{theorembox}{MAP Equivalence to Regularized Machine Learning}
\index{Regularized Machine Learning}\index{Ridge Regularization}\index{Lasso Regularization}
Let $Y_i = \mathbf{w}^T \mathbf{x}_i + \epsilon_i$ with Gaussian noise $\epsilon_i \sim \mathcal{N}(0, \sigma^2)$. The log-likelihood is $-\frac{1}{2\sigma^2} \sum (y_i - \mathbf{w}^T \mathbf{x}_i)^2$.
\begin{itemize}
    \item \textbf{Gaussian Prior $\mathbf{w} \sim \mathcal{N}(\mathbf{0}, \tau^2 \mathbf{I})$ (Ridge / $L_2$ Regularization):}
    $$\ln \pi(\mathbf{w}) = -\frac{1}{2\tau^2} \|\mathbf{w}\|_2^2 + C \implies \hat{\mathbf{w}}_{\text{MAP}} = \arg\min_{\mathbf{w}} \left\{ \sum_{i=1}^n (y_i - \mathbf{w}^T \mathbf{x}_i)^2 + \frac{\sigma^2}{\tau^2} \|\mathbf{w}\|_2^2 \right\}$$
    \item \textbf{Laplace Prior $\pi(w_j) = \frac{\lambda}{2} e^{-\lambda |w_j|}$ (Lasso / $L_1$ Regularization):}
    $$\ln \pi(\mathbf{w}) = -\lambda \|\mathbf{w}\|_1 + C \implies \hat{\mathbf{w}}_{\text{MAP}} = \arg\min_{\mathbf{w}} \left\{ \sum_{i=1}^n (y_i - \mathbf{w}^T \mathbf{x}_i)^2 + 2\sigma^2 \lambda \|\mathbf{w}\|_1 \right\}$$
\end{itemize}
\end{theorembox}

\begin{videocard}{IITM BS Lecture: Bayesian Estimation: Finding Estimators}{https://www.youtube.com/watch?v=2owMxamj-Tc}
Official lecture solving MAP and Bayes Mean estimators for Beta-Binomial and Normal conjugate models.
\end{videocard}

\section{Exam Gotchas \& Cognitive Traps}

\begin{trapbox}{MAP Estimator is NOT Invariant Under Reparameterization}
\textbf{The Trap:} Assuming that if $\hat{\theta}_{\text{MAP}}$ is the MAP estimator of $\theta$, then $g(\hat{\theta}_{\text{MAP}})$ is the MAP estimator of $g(\theta)$.

\textbf{The Reality:} Unlike the MLE (which is strictly invariant), the \textbf{MAP estimator is NOT invariant} under nonlinear transformations!
Because the posterior is a probability \emph{density}, transforming variables introduces a Jacobian factor:
$$f_\Phi(\phi \mid \mathbf{x}) = f_\Theta(g^{-1}(\phi) \mid \mathbf{x}) \cdot \left| \frac{d\theta}{d\phi} \right|$$
This derivative factor shifts the peak (mode) of the density surface! Consequently:
$$\hat{\phi}_{\text{MAP}} \neq g(\hat{\theta}_{\text{MAP}})$$
\end{trapbox}

\begin{trapbox}{Credible Intervals vs. Confidence Intervals}
\textbf{The Trap:} Confusing a Bayesian $95\%$ Credible Interval with a Frequentist $95\%$ Confidence Interval.

\textbf{The Reality:}
\begin{itemize}
    \item \textbf{Frequentist Confidence Interval:} The parameter $\theta$ is a fixed constant. You cannot assign a probability to $\theta$. The $95\%$ probability applies solely to the algorithmic coverage of hypothetical future datasets.
    \item \textbf{Bayesian Credible Interval:} The parameter $\theta$ is an active random variable with a posterior density $f(\theta \mid \mathbf{x})$. A $95\%$ credible interval $[L, U]$ satisfies $\int_L^U f(\theta \mid \mathbf{x}) d\theta = 0.95$. Here, you \textbf{can} legitimately state: "Given the observed sample data and prior, there is a $95\%$ probability that the true parameter $\theta$ lies within $[L, U]$!"
\end{itemize}
\end{trapbox}

\section{Comparative Synthesis Matrix}

\begin{table}[htbp]
\centering
\small
\begin{tabular}{p{3.2cm}p{3.8cm}p{3.8cm}p{3.5cm}}
\toprule
\textbf{Criterion} & \textbf{Frequentist MLE} & \textbf{Bayesian MAP} & \textbf{Bayesian Posterior Mean} \\
\midrule
\textbf{Mathematical Target} & $\arg\max_\theta L(\theta)$ & $\arg\max_\theta [L(\theta) \pi(\theta)]$ & $\E[\theta \mid \mathbf{x}] = \int \theta f(\theta|\mathbf{x}) d\theta$ \\
\addlinespace
\textbf{Loss Function} & Asymptotic likelihood optimization & 0--1 Loss (Infinitesimal hit-or-miss) & Squared Error Loss $(\theta - \hat{\theta})^2$ \\
\addlinespace
\textbf{Invariance Property} & \textbf{Strictly Invariant}: $\widehat{g(\theta)} = g(\hat{\theta})$ & \textbf{Not Invariant} (Jacobian shifts mode) & Not invariant ($\E[g(\theta)] \neq g(\E[\theta])$) \\
\addlinespace
\textbf{Small-Sample Risk} & High variance, prone to zero-count overfitting & Regulated by prior strength (shrinkage) & Optimal under MSE, incorporates full curve \\
\addlinespace
\textbf{Machine Learning Analog} & Unregularized OLS, Logistic Regression & Ridge ($L_2$) and Lasso ($L_1$) regularized regression & Ensemble model averaging, Bayesian Neural Nets \\
\bottomrule
\end{tabular}
\caption{Comparison of Estimation Philosophy, Properties, and Decision Loss Functions.}
\end{table}

\newpage
\section{Practice Exercises}
\textit{Detailed step-by-step solutions are available in the Appendix.}

\subsection*{Part A: Analytical \& Computational Problems}

\begin{enumerate}
\item \textbf{Beta-Binomial Conjugate Posterior:} Let coin flip success probability $p \sim \text{Beta}(\alpha = 2, \beta = 2)$ with prior PDF $f(p) = 6p(1-p)$ for $0 \le p \le 1$. In $n = 5$ coin flips, $k = 4$ heads are observed ($L(p) = \binom{5}{4} p^4 (1-p)^1$).
    \begin{enumerate}
        \item Find the posterior distribution $f(p \mid k = 4)$.
        \item Compute the MAP estimator $\hat{p}_{\text{MAP}}$.
        \item Compute the Bayes posterior mean estimator $\hat{p}_{\text{Bayes}}$.
    \end{enumerate}

\item \textbf{Normal-Normal Conjugate Model:} Let observation $X \mid \mu \sim \mathcal{N}(\mu, \sigma^2 = 4)$. A prior is placed on $\mu \sim \mathcal{N}(\mu_0 = 0, \sigma_0^2 = 1)$. Single observation $x = 5$ is made. Compute the posterior mean $\E[\mu \mid X = 5]$.

\item \textbf{Uniform Prior MAP vs MLE Comparison:} Let $X \sim \text{Binomial}(n = 10, p)$. If we assume a uniform prior $p \sim \text{Uniform}(0, 1)$, show that $\hat{p}_{\text{MAP}} = \hat{p}_{\text{MLE}} = \frac{x}{10}$.

\item \textbf{Laplace Prior and $L_1$ Regularization Equivalence:} Show that maximizing the log-posterior $\ln f(\mathbf{w} \mid \mathbf{x}, y) = \ln f(y \mid \mathbf{x}, \mathbf{w}) + \ln f(\mathbf{w})$ under a zero-mean Laplace prior $f(w_j) = \frac{\lambda}{2} e^{-\lambda |w_j|}$ adds an $L_1$ penalty term $-\lambda \|\mathbf{w}\|_1$ to the log-likelihood (Lasso regression).

\item \textbf{Data Science Application (Bayesian AB Testing Click Probability):} In an email campaign, variant A conversion probability $p$ is assigned prior $\text{Beta}(1, 1)$ (Uniform). After sending $n = 100$ emails, $k = 15$ conversions occur. Compute the posterior distribution $\text{Beta}(\alpha', \beta')$ and the posterior mean estimate for conversion rate.
\end{enumerate}

\subsection*{Part B: Conceptual Diagnostic Traps (Exam Gotchas)}

\begin{enumerate}[resume]
\item \textbf{Diagnostic Trap (Loss Function and Bayes Estimator Selection):}
A machine learning engineer constructs a continuous posterior distribution $f(\theta \mid \mathbf{x})$. Under which of the following loss functions is the optimal Bayes point estimator mathematically guaranteed to be the posterior mean $\E[\theta \mid \mathbf{x}]$?
\begin{enumerate}
    \item Absolute Error Loss: $L(\theta, \hat{\theta}) = |\theta - \hat{\theta}|$
    \item Zero-One Loss: $L(\theta, \hat{\theta}) = \mathbb{I}(|\theta - \hat{\theta}| > \epsilon)$ as $\epsilon \to 0$
    \item Squared Error Loss: $L(\theta, \hat{\theta}) = (\theta - \hat{\theta})^2$
    \item Exponential Loss: $L(\theta, \hat{\theta}) = e^{c(\theta - \hat{\theta})}$
\end{enumerate}

\item \textbf{Diagnostic Trap (Reparameterization Invariance):}
Suppose $\hat{\theta}_{\text{MLE}}$ is the MLE of parameter $\theta$, and $\hat{\theta}_{\text{MAP}}$ is the MAP estimator of $\theta$ under a continuous non-constant prior $\pi(\theta)$. Let $g(\theta) = \theta^3$. Which of the following statements is strictly correct?
\begin{enumerate}
    \item Both the MLE and MAP estimators satisfy invariance: $\widehat{g(\theta)}_{\text{MLE}} = (\hat{\theta}_{\text{MLE}})^3$ and $\widehat{g(\theta)}_{\text{MAP}} = (\hat{\theta}_{\text{MAP}})^3$.
    \item Neither the MLE nor MAP estimator is invariant under non-linear transformations.
    \item The MLE is invariant ($\widehat{g(\theta)}_{\text{MLE}} = (\hat{\theta}_{\text{MLE}})^3$), but the MAP estimator is generally not invariant due to the Jacobian of transformation.
    \item The MAP estimator is invariant, but the MLE fails invariance whenever the prior is non-uniform.
\end{enumerate}
\end{enumerate}
